% =====================================================================================
%  Discussion, threats to validity, limitations. Included by paper.tex after Results.
% =====================================================================================

\section{Discussion}

\subsection{A fifth of the signal was not there}

The headline cost figure is \num{0.0792} AUC, or \SI{19.6}{\percent} of everything the
unconstrained model achieved above chance. That is not a rounding correction. A practitioner
deploying the unconstrained model would find a fifth of its apparent discrimination absent
at run time, because the attributes carrying it are recorded only after the decision has been
taken. The same measurement under an \num{80}/\num{20} split differs by less than
\num{0.002} AUC, so the estimate is a property of the data rather than of the split.

This is consistent in direction with reported effects elsewhere, including the
\num{15.07} point accuracy drop \citet{mishra2026taxonomy} measure on build prediction and
the example-leakage rates above \SI{80}{\percent} that \citet{abb2024discussion} find on
standard process logs. The contribution here is that the cost and the benefit are measured on
the same log under one protocol, so neither can be quoted without the other.

\subsection{Admissibility is not the same as parsimony}

A natural response to a leakage finding is to use fewer attributes. The ladder shows that
would be the wrong lesson. Discrimination rises monotonically from \num{0.5100} at three
admissible attributes to \num{0.8270} at seventeen, a gain of \num{0.2560} AUC, reproduced to
within \num{0.0053} across split ratios. Restricting to a small attribute set costs far more
than enforcing admissibility does.

The distinction matters because the two are easily conflated. Attributes should be dropped
because they are unavailable at decision time, not because they are numerous. Nine attributes
are excluded here, each with its reason recorded; the remaining seventeen are kept because
each is observable when the assignment is made.

\subsection{Choosing the admissible target can improve accuracy, not only defensibility}

The service level result is the one we did not expect. Breach cannot be modelled directly as
a survival event on this log without the model reading its own label, because breach is
defined by duration exceeding a per-priority threshold and the measured breach rate past that
threshold is exactly \num{1.0000}. An early version scored \num{0.9418} AUC on precisely that
artifact.

Modelling resolution instead, and deriving breach as the conditional survival ratio
$S(\tau \mid x) / S(t \mid x)$, produces \num{0.8373} AUC at intake against \num{0.7879} for a
direct classifier. The admissible formulation is better by \num{0.0494} AUC. Admissibility is
usually framed as a constraint that costs performance. Here, choosing the object that can
actually be estimated improved it, because the resolution hazard is a genuine random quantity
given elapsed time while breach, past the threshold, is not.

\subsection{The weighting objection answers itself}

Objective weighting methods are known to disagree, to be unstable under perturbation, and to
admit rank reversal \citep{mukhametzyanov2021specific,linh2025evaluating,li2026trend}. CRITIC
is additionally unsuitable in isolation here, because two of the seven criteria are model
outputs whose spread reflects calibration rather than importance, and isotonic calibration
compresses variance. Reporting CRITIC alone would convert a modelling artifact into a weight
and label it data-driven.

Reporting three schemes resolves this empirically rather than rhetorically. The induced
assignment orderings agree at Spearman $\rho = \num{0.959}$ between CRITIC and uniform
weighting and $\rho = \num{0.946}$ between CRITIC and entropy. The weighting choice is
therefore not load-bearing for these results, which is a stronger statement than defending any
one scheme.

\subsection{The modifications work, but the base algorithm matters more}

The proposed operator set does improve on standard NSGA-II, significantly and on four of five
metrics (Section~\ref{sec:assignment-results}). It is then outperformed by NSGA-III, an
off-the-shelf control included only as a many-objective reference, which reaches higher
hypervolume from a front roughly a third the size. Two considerations make this informative
rather than disappointing.

First, five objectives place this problem in the many-objective regime where dominance-based
selection is known to degrade, a limitation characterised when NSGA-III was introduced
\citep{debjain2014evolutionary} and still under active theoretical study
\citep{zheng2024runtime,doerr2026difficulties}. A result in which a reference-point method
outperforms a dominance-based one on five objectives agrees with that theory.

Second, \citet{wang2025rethinking} argue that much claimed novelty in metaheuristic research
does not survive controlled comparison. An ablation that isolates each proposed modification
is the appropriate response to that argument, and it is why the ablation is reported per
component rather than only in aggregate.

\subsection{What the cost looks like in cases rather than in AUC}

A \num{0.0792} AUC loss is an abstraction. The confusion matrices in
Section~\ref{sec:confusion} put it in units an operations manager can act on. On the same
\num{13982} held-out incidents, the unconstrained model catches \num{4295} of the \num{5136}
breaches and the admissible one catches \num{3795}. The \num{500} cases in that difference are
not a modelling subtlety. They are the incidents an escalation queue built on the
unconstrained model would appear to catch and would not.

The direction of the error also changes. The unconstrained model is roughly symmetric,
\SI{82.3}{\percent} specificity against \SI{83.6}{\percent} recall. The admissible model is
not: \SI{74.6}{\percent} specificity against \SI{73.9}{\percent} recall. Removing the
post-hoc attributes removes signal about the long-running cases specifically, because those
attributes counted the events a long-running case accumulates. What is left behind
discriminates the ordinary case well and the pathological case poorly, which is the harder
half of the problem and the half a service desk actually cares about.

\subsection{The model is data-limited, not capacity-limited}

The learning curve in Section~\ref{sec:learning} separates two explanations for
\num{0.8270} that are often conflated. Training AUC is \num{0.9999} at \num{1631} incidents,
so the learner has ample capacity; the ceiling is not the model class. The generalisation gap
falls monotonically from \num{0.2234} to \num{0.0563} across a twentyfold increase in data
and is still falling, and held-out AUC is still rising at the largest training size available.

Two consequences follow, and the second is a warning against the obvious fix. More incidents
from this desk would help, so the practical route to a better admissible model is more history
rather than a better algorithm, which is consistent with the four learners landing within
\num{0.0158} AUC of each other. And heavier regularisation is the wrong first response: it
would close the gap by moving the training curve down, not by moving the held-out curve up.
The gap is not evidence that the model is over-parameterised. It is evidence that
\num{32623} incidents is not yet enough for a 17-attribute problem with this much label noise.

\subsection{Ranking and operating point fail separately}

The chronological result is usually reported as one number, and Section~\ref{sec:temporal}
shows it is two failures with different remedies. Ranking degrades from \num{0.8258} to
\num{0.8074} AUC and no threshold recovers that. The operating point degrades as well, from
$F_1 = \num{0.6790}$ on the random split to $F_1 = \num{0.6087}$ at the same default cut point,
and tuning the cut point on training-period data alone recovers it to \num{0.6157} without
retraining anything. Recall is where the operating point actually moves: \SI{67.7}{\percent} at the
default against \SI{81.0}{\percent} tuned.

Conflating the two produces the wrong operational conclusion in both directions. A team that
sees only the fall in $F_1$ will retrain a model that did not need retraining. A team that
sees only the tuned $F_1$ will not learn that its threshold needs re-tuning as the process
drifts, which on this log it does within a year: median handle time falls from \num{4.28} to
\num{2.79} hours between the training and evaluation periods.

\subsection{Six defects in our own work, and what each one cost}
\label{sec:harness-defects}

Six errors were found in this work and corrected before any number reported here was
produced. All six are recorded because a paper about claims that do not survive verification
should hold its own instrument to the same standard, and because an earlier version of this
section did not: it named three and gave the other three either a sentence elsewhere or
nothing.

That omission was itself the more serious error, so the ordering here is deliberate. The
first three below are defects in what was measured. They change what the paper is entitled
to claim. The last three are defects in the instrument. They change how reliably a number
can be reproduced, which matters, but a reproducible measurement of the wrong thing is
still the wrong thing.

\paragraph{Validity 1: the incident file was a partial copy.} The working copy of
\texttt{Detail\_Incident.csv} held \num{31238} rows against the canonical
\num{46606}, a third of the log missing, and it was a prefix rather than a sample. The file is
ordered by incident ID, which tracks open time, so a prefix is close to a truncation in time.
That is why it cost so much more on the chronological arm than on the random one:
\num{0.0307} against \num{0.0043}, measured by reintroducing it alone. The canonical file was
re-downloaded from the 4TU DOI, and the digest of every raw input is now recorded in
\texttt{code/input\_checksums.json} and checked by \texttt{evaluation/verify\_input\_data.py};
the incident header is MD5 \texttt{b0608464312c135da289d99911735a5d}. A digest stored on the first day
would have caught this on the first run, and every artifact computed from the partial file
was computed correctly from the wrong input, which is the failure mode a checksum exists
for.

\paragraph{Validity 2: a feature of ours violated our own Rule 1.} The standing-queue
count at a case's arrival was computed over cases that eventually resolved, excluding the
\num{1780} that never do. Whether a case will ever resolve is not knowable at
$t^{\mathrm{dec}}$, so the feature read the future. It is worth stating plainly that this
is a violation of the rule this paper exists to state, committed by its authors, and that it
was found by applying the rule mechanically rather than by intuition. Reintroducing it alone
raises the admissible AUC by \num{0.0015}, which is the direction a leak should move a
score and is the reason we believe the diagnosis.

\paragraph{Validity 3: Rule 2 was not implemented as Rule 2.} The pipeline refitted the
group encoding inside every cross-validation fold. That prevents an encoding from seeing the
rows in its own evaluation fold, and it is what most careful work in this area does, but it
is not Equation~\ref{eq:rule2}: under a random split a case from 2012 could still be encoded
using cases that resolved in 2014. The paper's contribution is that equation, so for a period
the protocol described was not the protocol that ran. It is now implemented as written, in
\texttt{code/causal\_encoding.py}, checked against an $O(n^2)$ reference implementation to a
maximum absolute difference of zero. Reintroducing the fold refit alone raises the admissible
AUC by \num{0.0043}. The distinction is therefore load-bearing for the claim and not for the
number, and both halves of that sentence are reported.

\paragraph{Instrument 1: the seed was not a seed.} It was derived using Python's built-in string hash, which is salted per process. The
same variant therefore drew a different seed on every invocation, measured at \num{43},
\num{41} and \num{32} across three processes for one variant name. The harness documentation
asserted reproducibility that did not hold. It now uses a stable checksum, and two independent
invocations produce bitwise-identical metrics.

\paragraph{Instrument 2: the hypervolume reference point moved.} It was taken from whichever run completed first, making every
hypervolume value dependent on variant ordering and incomparable between runs with different
variant sets. It is now the nadir across all variants and all runs on a batch, computed before
any hypervolume is taken.

\paragraph{Instrument 3: the thread count was never pinned.} This one most directly
threatens what the paper claims about verification. Every number here is
printed to four decimals and checked against a tolerance of $\num{0.00005}$. That tolerance
only means something if a re-run reproduces the number, and it did not. Running the headline
pipeline with one seed, one dataset and one split, changing nothing but the OpenMP thread
count, moves the admissible AUC by \num{0.0015} (2 gives \num{0.8267}, 4 gives \num{0.8258}, 8 gives \num{0.8273})
and the unconstrained AUC by \num{0.0003}. That is roughly \num{30} times the tolerance,
from a variable that appeared in no configuration file and in no methods section.

The cause is not randomness and no seed can absorb it. Gradient-boosted implementations
reduce per-thread histogram partials in thread order, floating-point addition is not
associative, and the resulting difference in a split point compounds through the trees. The
control confirms it: \num{2} runs at the same thread count are bitwise identical,
so the variation is entirely explained by how the work was divided, not by what the code did
with it. The thread count is now pinned in the harness and recorded in every run log, and the
measurement above is reproduced by \texttt{evaluation/thread\_determinism.py}.

We report it rather than quietly pinning it because the implication generalises past this
paper. A reviewer who re-runs a published gradient-boosted pipeline on different hardware
should not expect the fourth decimal to agree, and a paper that verifies its own numbers to
four decimals without fixing the thread count is verifying them against a moving target. Ours
was. Pinning does not make the numbers portable across CPU architectures either, since vector
width changes the summation order within a thread as well; it removes the one source that
varied on a single machine, which was the source that had been silently moving our own
results between runs.

The first two instrument defects do not affect the predictive results in
Sections~\ref{sec:cost} to \ref{sec:hazard-results}, which come from a separate pipeline. Both
invalidated every earlier assignment run, and those were discarded.

The consequence is worth stating plainly rather than in passing. Under the uncorrected
harness the proposed method appeared to LOSE to standard NSGA-II on hypervolume. Corrected,
it wins, significantly. The defect did not merely add noise: it reversed the sign of the
study's headline comparison. Had the harness not been audited, this paper would have
reported the opposite conclusion with the same apparent statistical confidence.

\subsection{Threats to validity}

\textbf{Construct.} The service level label is derived from handle time against \emph{twice} the
per-priority median rather than from a contractual target, because the log's own
\texttt{Alert Status} field is uniformly \texttt{closed} across \num{46606} of \num{46809} raw
incidents and carries no usable breach signal. Results should be read as conditional on that
operationalisation, and Equation~\ref{eq:label} states it exactly so it can be varied. On the
random-split arms the median is taken over the whole log. That is a property of the label and
not of any feature, no model observes it, and every configuration is scored against the same
fixed label, so it cannot affect the comparison between configurations; the chronological arm
re-derives it from training-period cases alone, where it would otherwise matter.

\textbf{Statistical.} The fold-level Wilcoxon signed-rank test at $n=5$ has a smallest
attainable two-sided $p$-value of \num{0.0625} and therefore cannot reject at
$\alpha=\num{0.05}$ regardless of effect size. Every fold-level comparison in
Section~\ref{sec:cv} returns exactly that floor, which means one learner won on all five folds
and not that the difference is small. Model comparisons are therefore made with 2000-resample
bootstrap intervals on the test set, and the fold-level result is reported for transparency
rather than used as evidence.

\textbf{Internal.} Every parameter is fitted on the earliest \SI{80}{\percent} of the timeline,
and the escalation cut point is tuned on training-period data alone. The chronological holdout
is reported alongside the random split precisely because the gap between them, \num{0.8270}
against \num{0.8085}, is itself informative about drift. That gap was \num{0.1326} before the date-parsing defect in the loader was corrected and is \num{0.0185} after it, so most of what earlier drafts read as distribution shift was the shift in how their own timestamps were parsed.

\textbf{External.} The protocol is measured on five public event logs spanning four domains:
IT service management, consumer lending, acute care, healthcare administration and public
administration. That is enough to say the protocol transfers and enough to say the magnitude
does not. The share of above-chance signal lost to admissibility ranges from \SI{9.7}{\percent}
on Hospital Billing to \SI{69.5}{\percent} on Road Traffic Fines. Any single figure quoted from
one log, including the one this study began with, is a point on that range and not a constant.

Four of the five logs are Dutch or Italian European public releases and all five are historical.
Nothing here speaks to a live deployment, to a non-European process, or to an organisation that
knows its assignment policy is being measured.

\textbf{Conclusion.} Statistical comparisons use paired Wilcoxon signed-rank tests with
Holm-Bonferroni correction across five metrics, and every run is seeded and reproducible.

\section{Limitations}

\textbf{An IT service desk is not a sales pipeline.} This work is motivated by lead
assignment, and the objective analogous to conversion is first-time-right resolution. No
public dataset carries real leads, real individual representative identifiers and real
conversion outcomes together. That absence is a finding rather than an oversight, and it is
why the evaluation is conducted on an operational log whose assignment decisions and outcomes
are real even though its domain is not sales.

\textbf{A parsing defect in an earlier version of this work, and what it cost.} The incident
file records timestamps day-first; an earlier version of this pipeline parsed them month-first,
which silently transposed day and month on the \SI{41.0}{\percent} of rows where both fields
are at most twelve. Two consequences were reported as properties of the data and were not.
The observation window was given as \num{1055} days against a true \num{785}. And the
first-assignment delay appeared corrupted in its upper tail, with a median of \num{16.2}
hours against a 90th percentile of \num{5684}, which led to \SI{35.6}{\percent} of rows being
excluded from the freshness and urgency fits as log artifacts. Under the correct parse that
attribute has a median of \num{0.34} hours and a 90th percentile of \num{15.48} over the
\num{46605} incidents of the predictive matrix, and the exclusion above the \num{168}-hour cap
falls to \SI{0.4}{\percent}. Table~\ref{tab:fitted} quotes \num{0.49} and \num{16.57} for the same
attribute because the formulations are fitted on the \num{39449} incidents of the assignment
working set rather than on all of them; the two populations differ and so do their medians. Nothing was corrupt. The subtraction was mixing two
date conventions. Every number in this paper is computed after that correction, and the
defect is recorded here because a paper about attributes that mislead should say when its
own instrument did.

\paragraph{What the corrections do not account for.} Each of the six was reintroduced on
its own, everything else held corrected, and the pipeline re-run
(\texttt{evaluation/attribution\_run.py}). The date defect explains most of the change on
the random split, \num{0.0389} of \num{0.0399}, and the truncated file explains most of it
on the chronological arm, \num{0.0307} of \num{0.1543}. But all six reintroduced together
give \num{0.7996} and \num{0.7235}, while the figures an earlier version of this work
published were \num{0.7859} and \num{0.6533}. So the reconstruction accounts for about two
thirds of the change on the random split and just over half of it on the chronological arm.
The remaining \num{0.0137} and \num{0.0702} are \emph{unexplained}. We do not attribute
them to the thread count: that effect is measured above at \num{0.0015} and \num{0.0019},
which is forty times too small to cover the chronological residual, and offering it as a
candidate would be exactly the kind of unfalsifiable explanation this paper objects to
elsewhere. The earlier pipeline was rewritten rather than parameterised, so it cannot be
re-run to settle the question, and a reader is entitled to treat the pre-correction figures
as unreproducible rather than as merely superseded.

\paragraph{Whether the corrections are self-serving.} They moved the headline from
\num{0.7859} to \num{0.8258} and the chronological arm from \num{0.6533} to \num{0.8076}, and a
round of bug fixing that improves nearly every number deserves the suspicion it attracts. Two
things answer it, and neither is an assurance. First, the two leaks we removed moved the score
\emph{up} when reintroduced, \num{0.0015} and \num{0.0043}, which is the direction a leak must move it;
had removing a leak raised our score, the diagnosis would have been wrong. Second, and more
to the point, several corrections cost us. The four learners are no longer statistically
distinguishable, where an earlier draft reported two of three gaps excluding zero. Two of
the three proposed algorithm components now \emph{raise} the objective when removed, so
they contribute nothing. A component described as paying for itself twice in fact costs
\SI{77}{\percent} more runtime, and that claim is withdrawn. And the study's most quotable finding,
that a deployed model would recover under a tenth of future breaches at its default
threshold, was an artifact of the date defect: the true figure is \SI{67.7}{\percent}, and the
dramatic version is withdrawn rather than softened. A correction round that only ever helped
would be evidence of selection. This one did not.

\textbf{Counterfactual outcomes are not observed.} A historical log records only the
assignment that happened. Any policy differing from the recorded one is scored by a model
rather than by ground truth, which is the standard objection to optimisation studies on
logged data. We report the fraction of decisions matching the recorded assignment together
with the real observed breach rate on exactly those cases, as the only non-simulated outcome
signal available. Off-policy evaluation would address this properly and is left to future
work.

\textbf{Group-level and individual-level assignment both appear, and the distinction matters.}
BPI 2014 identifies assignment groups rather than individuals, so representative-level features
such as experience are not estimable on it. Three of the four added logs do carry individual
resource identifiers, and the encoding runs at that level there: \num{120} distinct originators
on BPI Challenge 2017 and \num{894} on Hospital Billing. The protocol is unchanged between the
two settings, because Rule 2 is stated over whatever entity the log identifies. What we cannot
say is whether encoding at individual level is better than at team level, since no log here
offers both and a comparison across logs would confound the entity with everything else.

\textbf{Explanation is reported but not validated against human decisions.} SHAP attributions
are computed for the predictive layer, but attribution methods have known impossibility
results \citep{bilodeau2024impossibility} and sensitivity to feature representation
\citep{hwang2025shapbased}, and whether explanations improve human decisions is unsettled
\citep{haag2026explanations,chae2025explanation}. No claim is made that these explanations
improve operator performance.
